% Methods text: calibration of measurement discrepancies + aggregation of asynchronous recordings (21S).
% Draft 2026-11-10; numbers from results/reliability_model_v1.json and results/dc_reconciled_v2.pkl (PlexusFlow21S).
% Checked against src/plexusflow/analyze/profile_flux.py (parabolicity), src/plexusflow/reconcile/reliability.py
% (c model, fit) and src/plexusflow/reconcile/dc.py (vessel_input, reconcile_dc).

\subsection*{Calibration of measurement discrepancies}

For each vessel recording we characterised measurement quality by two properties of the time-averaged velocity field:
its agreement with a parabolic flow profile and its alignment with the vessel axis. Let $k$ index vessels and $t$ the
image tiles (recordings) in which vessel $k$ was measured. Flux was estimated from cross-sections sampled densely along
the vessel centreline; at each cross-section $i$ a parabolic profile was fitted to the axial velocity. We defined the
mean parabolicity

\begin{equation}
p_{k,t}
=
\frac{1}{N_{k,t}}
\sum_{i=1}^{N_{k,t}}
\operatorname{clip}\!\left(R_{k,t,i}^2,0,1\right),
\label{eq:parabolicity}
\end{equation}

where $N_{k,t}$ is the number of cross-sections, $R_{k,t,i}^2$ is the coefficient of determination of the parabolic fit
at cross-section $i$, and cross-sections without a valid fit (too few velocity samples, implausible fitted radius, or no
central velocity maximum) contribute zero. Unlike the acceptance criterion used to estimate flux, which retains only
cross-sections with $R^2 \geq 0.5$, this continuous descriptor summarises profile quality along the whole vessel.
(A recording in which fewer than two accepted cross-sections agreed in flow direction gave no flux and the vessel was
treated as unmeasured in that tile.) We
also computed the directional alignment $a_{k,t}$, the speed-weighted mean of $|\cos\vartheta|$, where $\vartheta$ is the
angle between the time-averaged velocity vector and the local centreline tangent, over pixels within one vessel radius
of the centreline.

We combined these descriptors into a dimensionless fractional measurement-discrepancy scale,

\begin{equation}
c_{k,t}
=
c_0
+ c_p(1-p_{k,t})
+ c_a(1-a_{k,t}),
\label{eq:measurement_uncertainty}
\end{equation}

where $c_0$ is the baseline discrepancy scale and $c_p$ and $c_a$ describe its dependence on profile quality and flow
alignment. Smaller $c_{k,t}$ indicates a more reliable recording. Under the working error model, $c_{k,t}$ specifies the
modelled fractional standard deviation of measurement discrepancies, apart from an additive floor: a recording with $c_{k,t}=0.3$ is assigned a discrepancy standard deviation of about $30\%$ of its flux. This is a modelled
scale, not a bound on the error.

We calibrated the coefficients using conservation at junctions whose incident vessels were all measured
simultaneously, within a single image tile $t$. For such a junction $j$, the measured net flux is

\begin{equation}
r_j
=
\sum_{k\in\mathcal E_j}s_{jk}\,y_{k,t},
\end{equation}

where $y_{k,t}$ is the measured mean flux and $s_{jk}=\pm1$ specifies its orientation relative to the junction. For fully
resolved junctions, conservation requires the underlying mean fluxes to balance. The observed imbalance therefore provides
a measure of the combined discrepancies of the incident flux estimates, without requiring the individual true fluxes to
be known. We modelled the junction imbalances as zero-mean Gaussian
with variance

\begin{equation}
V_j(\boldsymbol\theta)
=
\sum_{k\in\mathcal E_j}
c_{k,t}^{\,2}\,y_{k,t}^{\,2}
+\sigma_0^2,
\label{eq:junction_error_variance}
\end{equation}

where $\sigma_0$ is an absolute junction-discrepancy floor and $\boldsymbol\theta=(c_0,c_p,c_a,\sigma_0)$. This
formulation assumes that the vessel-level discrepancies are approximately uncorrelated, with variances proportional to
squared flux; the measured fluxes stand in for the unknown underlying fluxes when evaluating the variance. The
parameters were estimated by minimising the Gaussian composite negative log-likelihood,

\begin{equation}
\hat{\boldsymbol\theta}
=
\arg\min_{\boldsymbol\theta}
\frac12\sum_j
\left[
\log V_j(\boldsymbol\theta)
+\frac{r_j^2}{V_j(\boldsymbol\theta)}
\right],
\label{eq:uncertainty_calibration_likelihood}
\end{equation}

with all four parameters constrained to be positive (optimised on a log scale). The likelihood is a composite one:
neighbouring junctions can share a vessel measurement, so their imbalances are not all independent. Fitted to all
$4{,}411$ single-tile junctions, the model gives $\hat c_0=0.054$, $\hat c_p=0.42$, $\hat c_a=1.45$ and $\hat\sigma_0=1.3\times10^{-3}$\,nL\,s$^{-1}$.

These full-data coefficients are used in all subsequent analyses. Calibration was evaluated separately on held-out
tiles: the tiles were divided into five folds, the model was refitted on the junctions of four folds, and the
standardised imbalances $z_j=r_j/\sqrt{V_j}$ of the fifth were computed with these fold-specific coefficients. A calibrated
model gives a root-mean-square $z_j$ near unity at every level of measurement quality. We summarise the quality of a
junction by its effective fractional discrepancy,

\begin{equation}
c_{j,\mathrm{eff}}
=
\left(
\frac{\sum_{k\in\mathcal E_j}c_{k,t}^{\,2}\,y_{k,t}^{\,2}}
{\sum_{k\in\mathcal E_j}y_{k,t}^{\,2}}
\right)^{1/2},
\end{equation}

the predicted imbalance standard deviation as a fraction of the root-sum-square measured flux (ignoring $\sigma_0$), and
compare the quality-dependent model with a constant fractional-discrepancy model in which $c_{k,t}$ does not depend on
$p_{k,t}$ or $a_{k,t}$ (Fig.~\ref{fig:dc_reconciliation}a).

Calibration from junction conservation cannot identify a common multiplicative bias affecting all incident flux
measurements, such as a shared attenuation of measured velocities by depth averaging: scaling every incident flux by the
same factor leaves the balance unchanged. It likewise assumes approximately uncorrelated errors between the vessels of a
junction, and it does not account for physiological differences between tiles acquired at different times.

\subsection*{Aggregation of asynchronous measurements and conservative reconstruction}

Vessels in tile overlaps generally had several mean-flux measurements acquired at different times ($4{,}208$ of the
$7{,}120$ measured vessels in the 21-somite embryo). Rather than selecting a single recording, we formed a
quality-weighted mean for each vessel,

\begin{equation}
Q_k^{\mathrm{meas}}
=
\sum_t \alpha_{k,t}\,y_{k,t},
\qquad
\alpha_{k,t}=\frac{c_{k,t}^{-2}}{\sum_{t'}c_{k,t'}^{-2}},
\label{eq:measurement_aggregation}
\end{equation}

with all $y_{k,t}$ signed along a common orientation of vessel $k$. The quality-dependent component of the vessel-level
discrepancy scale is the weighted root-mean-square of the recording scales,

\begin{equation}
c_k^2
=
\sum_t \alpha_{k,t}\,c_{k,t}^2 .
\label{eq:aggregated_uncertainty}
\end{equation}

Because the recordings were acquired asynchronously, they cannot be treated as independent observations of a fixed
physiological state. We therefore retained a representative recording-level discrepancy scale rather than reducing $c_k$
according to the standard error of a weighted mean. We also computed the weighted between-recording spread,

\begin{equation}
s_{k,\mathrm{between}}^2
=
\sum_t \alpha_{k,t}\left(y_{k,t}-Q_k^{\mathrm{meas}}\right)^2,
\label{eq:between_recording}
\end{equation}

which was set to zero when only one recording was available. This term captures disagreement arising from measurement
variability, physiological change, or both; it is not an independently calibrated error variance, and adding it to the
quality-dependent term is a conservative modelling choice rather than a decomposition into independent sources.
Including it prevents strongly disagreeing recordings from producing an aggregated flux that is assigned excessive
weight, particularly when opposing measurements average to nearly zero.

The aggregated measurements are the inputs to the conservation-constrained reconstruction, with discrepancy scales

\begin{equation}
\sigma_k^2(Q'_k)
=
c_k^2\,(Q'_k)^2
+s_{k,\mathrm{between}}^2
+s_0^2,
\label{eq:reconstruction_uncertainty}
\end{equation}

where $Q'_k$ is the reconstructed reference flux and $s_0$ is a numerical floor, fixed at $1\%$ of the median absolute
measured flux ($2.2\times10^{-4}$\,nL\,s$^{-1}$ for the 21-somite embryo) and distinct from the fitted
junction-discrepancy floor $\sigma_0$ of Eq.~(\ref{eq:junction_error_variance}). Results were insensitive to its value:
with $s_0$ multiplied by $0.1$, by $5.8$ (equal to $\hat\sigma_0$) or by $10$, the total boundary inflow changed by at
most $0.4\%$ and the median per-vessel change of $Q'_k$ was at most $1\%$; the larger changes, including all changes of
direction, were confined to near-zero vessels. The scales $\sigma_k$ determine how strongly each aggregated measurement constrains
the reconstruction; they are not posterior uncertainties of the reconstructed fluxes. We therefore interpret the
reconstructed field as a conservation-consistent reference summarising asynchronous measurements, not as an
instantaneous flow field.

The reconstruction network was the vessel graph after removal of two kinds of annotated region, both drawn from the
videos rather than from the reconstructed field. Open regions, beside the dorsal aortae, are where blood enters the plexus
and where vessels are too dense and fast for reliable flux estimates; the vertices at which the remaining network meets
them, together with the two vitelline veins, form the boundary $\mathcal B$, through which flow may enter or leave.
Topology-uncertain regions are areas in which vessels exist but the graph does not represent them reliably (the bright
marginal band, sheet-like flow, and plexus too dense to trace). Each connected topology-uncertain region was replaced by a
mass-conserving compartment: the vertices on its boundary may exchange flow with one another, but the net exchange of the
region is zero, and its internal routing is not resolved. A topology-uncertain region contiguous with an open region was
treated as part of that open region. In the 21-somite embryo the reconstruction network has $8{,}969$ segments,
$38$ inflow vertices and three topology-uncertain regions ($575$ segments).

The reconstructed field was obtained by quasi-likelihood estimation under exact conservation. For each measured vessel
we defined the quasi-likelihood contribution through its derivative,

\begin{equation}
\frac{\partial \ell_k(q;Q_k^{\mathrm{meas}})}{\partial q}
=
\frac{q-Q_k^{\mathrm{meas}}}{c_k^2\,q^2+s_{k,\mathrm{between}}^2+s_0^2},
\qquad
\ell_k(Q_k^{\mathrm{meas}};Q_k^{\mathrm{meas}})=0,
\label{eq:quasi_likelihood}
\end{equation}

which integrates in closed form (a logarithm and an arctangent of $q$). Unlike weighted least squares with weights
fixed in advance, the variance is evaluated at the reconstructed flux rather than at the noisy measurement
\citep{Wedderburn1974}, so a vessel that happens to read low is not given spurious confidence. The reference field is

\begin{equation}
\mathbf Q'
=
\arg\min_{\mathbf q}
\;\frac1m\sum_{k\in\mathcal M}\ell_k(q_k;Q_k^{\mathrm{meas}})
+\frac{\varepsilon}{2}\sum_{k\notin\mathcal M}q_k^2
\quad\text{subject to}\quad
B_{\mathrm{int}}\,\mathbf q=\mathbf 0,
\label{eq:reconstruction}
\end{equation}

where $\mathcal M$ is the set of measured vessels, $B_{\mathrm{int}}$ is the signed incidence matrix restricted to
interior vertices, with one additional row per closed topology-uncertain region (the sum of the rows of its cut vertices),
$m$ is the median of
$1/\sigma_k^2$ at the starting point, and $\varepsilon=10^{-6}$ is a weak ridge on unmeasured vessels. The ridge fixes
circulations that conservation alone cannot determine (loops consisting only of unmeasured vessels); dividing by $m$
makes $\varepsilon$ relative to a typical measurement weight. Unmeasured vessels are therefore determined by
conservation, and the boundary fluxes $\mathbf f=B\mathbf Q'$ are outputs, not inputs.

Because the discrepancy variance depends on the reconstructed flux, the quasi-likelihood is generally nonconvex and may
admit multiple local minima. The starting point is therefore part of the estimator. We started from the uniform-weight conservative projection of the
measured fluxes and minimised Eq.~(\ref{eq:reconstruction}) by projected Newton iterations: each step is a weighted
conservative projection using a positive diagonal model of the Hessian, followed by Armijo backtracking, so the
objective decreases monotonically and the iteration stops at a constrained stationary point. To assess dependence on
the start, we repeated the solve from four perturbed starts (measured fluxes multiplied by $1+0.3\,\xi$,
$\xi\sim\mathcal N(0,1)$, fixed seeds). The median relative spread of $Q'_k$ across starts was $1.2\%$ for measured vessels (90th
percentile $19\%$) and $0.8\%$ for vessels with $c_k<0.5$. For $3.3\%$ of measured vessels, mostly low-flux ones, the
spread exceeded the reconstructed flux, so their direction is not determined by the data; together with unmeasured
vessels such cases carry $0.7\%$ of the summed absolute flux. Vessels whose direction, or whose magnitude by more than
$10\%$, changed between any two starts ($28\%$ of measured vessels) are flagged in the released data.
